\scanpage{149}
l'ensemble $\mathrm{Scind}\,\mathrm{isot}(E,F^\perp)$ des scindages isotropes de $(*)$ est stable \unsure{pour|par} \deleted{\unsure{B}} les opérations de $\mathrm{Alt}(G,G)\hookrightarrow\mathrm{Bil}(G,G)\simeq\mathrm{Hom}(G,F)\subset\mathrm{Hom}(G,F^\perp)$, et le quotient $\unsure{\mathrm{Sind}\,|\,\mathrm{Scind}\,}\mathrm{isot}/\mathrm{Alt}(G,G)$ est isomorphe (par la flèche canonique) à $\mathrm{Scind}/\underbrace{\mathrm{Bil}(G,G)}_{\substack{\shortparallel\\\mathrm{Hom}(G,F)}}$ (qui est lui-$=$ un torseur sous $\mathrm{Hom}(G,Q)$, où $Q=F^\perp/F$\ldots

\deleted{\unsure{d)}} e) Soit $\theta$ un scindage isotrope de $(*)$. Alors $(E,q)$ est somme directe quadratique de $F+\underbrace{\theta(G)}_{\substack{\shortparallel\\\check F}}$, \deleted{\unsure{avec sa}} \deleted{décomposé en \unsure{form}} \inserted{où $q$ est} sous forme décomposée $q(x,x')=\langle x,x'\rangle$, et de $Q_\theta$.

\deleted{\uline{Corollaire} Considérons un drapeau}
\[
\deleted{D:\quad0\subsetneq E_1\subsetneq E_2\subset\ldots\subsetneq E_{\unsure{r|\nu}}=F\qquad E.}
\]
\deleted{d'où un drapeau complété}
\[
\deleted{0\subsetneq E_1\ldots\subsetneq E_r=F\subset F^\perp\subsetneq E_{r-1}^{\perp}\subsetneq\ldots E_1^\perp\subsetneq F}
\]
\deleted{[\uline{NB} on peut avoir $F=F^\perp$ donc $Q=F^\perp/F=0$ \ldots]. Un scindage isotrope de $(*)$ définit alors une structure suppl. dans $E$ identifié à $F\oplus\check F\oplus Q$}
% archive sheet number: 148
